% ---- pg-395 （书页 383） ----

我们将不在这一点上解释因子. 像我们在 9.2 节注意过的那样, 因子(以及载荷)是唯一对正交旋转不变的. 因子旋转常常揭示一个简单的结构, 有助于进行解释. 在讨论了因子旋转之后我们将再次考虑这个例子(见例 9.9 和显示板 9.1). \hfill$\blacksquare$

\noindent\textbf{例 9.4}\quad (股票价格数据的因子分析)

例 8.5 介绍过 $p=5$ 支股票的 $n=103$ 周回报率组成的股票价格数据. 该例中, 从 $\mathbf{R}$ 得到过前两个样本主成分. 取 $m=1$ 和 $m=2$, 我们可以很容易地得到对正交因子模型的主成分分解. 特别的, 在以相应特征值的平方根调节之后, 被估计的因子载荷是样本主成分系数($\mathbf{R}$ 的特征向量). \hyperref[tab:9.2]{表 9.2} 列出估计的因子载荷、共性方差、特殊方差以及由各因子对 $m=1$ 和 $m=2$ 个因子解所解释的(标准化)样本总方差的比例. 共性方差由\hyperref[eq:9-17]{式(9-17)}给出. 所以, 例如在 $m=2$ 时, $\tilde{h}_1^2=\tilde{l}_{11}^2+\tilde{l}_{12}^2=0.732^2+(-0.437)^2=0.73$.

\begin{mvtable}
\phantomsection\label{tab:9.2}\mvtabcap{表 9.2}
\begin{tabular}{lccccc}
\toprule
\multirow{3}{*}{变量} & \multicolumn{2}{c}{单因子解} & \multicolumn{3}{c}{双因子解}\\
\cmidrule(lr){2-3}\cmidrule(lr){4-6}
& \makecell{因子载荷估计\\$F_1$} & \makecell{特殊方差\\$\tilde{\psi}_i=1-\tilde{h}_i^2$} & \multicolumn{2}{c}{因子载荷估计} & \makecell{特殊方差\\$\tilde{\psi}_i=1-\tilde{h}_i^2$}\\
\cmidrule(lr){4-5}
& & & $F_1$ & $F_2$ & \\
\midrule
1. J. P. Morgan & 0.732 & 0.46 & 0.732 & $-0.437$ & 0.27\\
2. Citibank & 0.831 & 0.31 & 0.831 & $-0.280$ & 0.23\\
3. Wells Fargo & 0.726 & 0.47 & 0.726 & $-0.374$ & 0.33\\
4. Royal Dutch Shell & 0.605 & 0.63 & 0.605 & 0.694 & 0.15\\
5. ExxonMobil & 0.563 & 0.68 & 0.563 & 0.719 & 0.17\\
\midrule
\makecell[l]{所解释的(标准化)\\样本总方差的\\累积比例} & 0.487 & & 0.487 & 0.769 & \\
\bottomrule
\end{tabular}
\end{mvtable}

对应于 $m=2$ 个因子的残差矩阵为
\[
\mathbf{R}-\tilde{\mathbf{L}}\tilde{\mathbf{L}}'-\tilde{\boldsymbol{\Psi}}=\begin{bmatrix}
0&-0.099&-0.185&-0.025&0.056\\
-0.099&0&-0.134&0.014&-0.054\\
-0.185&-0.134&0&0.003&0.006\\
-0.025&0.014&0.003&0&-0.156\\
0.056&-0.054&0.006&-0.156&0
\end{bmatrix}
\]
由双因子解所解释的总方差的比例要较单因子解大很多. 然而对 $m=2$, $\tilde{\mathbf{L}}\tilde{\mathbf{L}}'$ 得出的数一般说比样本相关要大, 对 $r_{13}$ 更是如此.

似乎很清楚了: 第一个因子 $F_1$ 代表总的经济条件, 可称之为\mvbold{市场因子}. 所有股票都高度地依赖于它, 且载荷都差不多相等. 第二因子将银行类股票和石油类股票进行对比(银行类在此因子上有比较大的负载荷而石油类有大的正载荷). 这样, $F_2$ 将不同行业股票加以区分, 可称之为\mvbold{行业因子}. 作为总结, 回报率要由总的市场条件和各个行业的独特运作以及残差或是公司特殊因子来共同决定, 这基本上是在例 8.5 中测试样本主成分所获得的结论. \hfill$\blacksquare$

\subsection{一个修正方法——主因子解}

有时考虑修改主成分方法. 我们用 $\mathbf{R}$ 的因子分析说明修改的理由, 虽然这个过程对 $\mathbf{S}$ 也是适用的. 如果因子模型 $\boldsymbol{\rho}=\mathbf{L}\mathbf{L}'+\boldsymbol{\Psi}$ 被准确指明了, 那这 $m$ 个公共因子将占据 $\boldsymbol{\rho}$ 的非对角元,
